% chapters/01_einleitung.tex
%
% Einleitung mit Forschungsfragen FF1 bis FF3 (Abschnitt 1.2).

\chapter{Einleitung}
\label{chap:einleitung}

Dieses Kapitel erläutert Motivation, Problemstellung und Zielsetzung der
vorliegenden Arbeit.
Abschnitt~\ref{sec:motivation_kontext} beschreibt den Anwendungskontext
und die beiden untersuchten Modellklassen.
Abschnitt~\ref{sec:problem_ziel} ordnet den Forschungsstand ein und leitet
daraus die Forschungsfragen und den Beitrag der Arbeit ab.
Abschnitt~\ref{sec:abgrenzung} legt den Untersuchungsumfang fest,
Abschnitt~\ref{sec:aufbau_arbeit} gibt einen Überblick über die weiteren
Kapitel.

\section{Motivation und Kontext}
\label{sec:motivation_kontext}

In Krankenhäusern, Nachrichtenredaktionen und Behörden entstehen täglich
zahlreiche Textdokumente, darunter Arztbriefe, Nachrichtenartikel und
Gerichtsentscheidungen.
Die darin enthaltenen Informationen liegen überwiegend unstrukturiert vor,
und ihre manuelle Erschließung ist aufwendig.
Natural Language Processing (NLP) umfasst Verfahren, mit denen solche
Texte automatisch verarbeitet und analysiert
werden~\cite{jurafsky2024speech,manning1999foundations}.
Eine grundlegende Aufgabe ist die Named Entity Recognition (NER):
Erwähnungen benannter Entitäten werden im Text lokalisiert und
vorgegebenen Kategorien zugeordnet~\cite{li2018survey}.
Die Ergebnisse dienen unter anderem als Grundlage für
Relationsextraktion, den Aufbau von Wissensgraphen und
Frage-Antwort-Systeme~\cite{jurafsky2024speech,keraghel2024recent}.

Welche Kategorien eine Anwendung benötigt, hängt vom Einsatzzweck ab.
Klassische Benchmarks wie CoNLL-2003 unterscheiden nur Personen,
Organisationen, Orte und eine Restkategorie~\cite{sang2003conll}.
In vielen Texten sind jedoch weitere Typen relevant, in Arztbriefen etwa
Krankheiten, in Nachrichtenartikeln Ereignisse oder Werke.
Feingranulare Datensätze wie MultiNERD bilden solche Typen ab und
unterscheiden 15 Kategorien, darunter Krankheiten, Ereignisse und
Medien~\cite{tedeschi2022multinerd}.

Die vorliegende Arbeit entstand am Deutschen Forschungszentrum für
Künstliche Intelligenz (DFKI) in der Forschungsgruppe Speech and Language
Technology (SLT).
Dort bildet die strukturierte Extraktion von Informationen aus
unstrukturierten Textbeständen ein zentrales Anwendungsfeld.

Für NER stehen heute zwei Klassen vortrainierter Sprachmodelle zur
Verfügung.
Encoder-Modelle wie BERT oder DeBERTa berechnen für jedes Token eine
kontextabhängige Repräsentation; ein darauf trainierter Klassifikator
weist jedem Token ein Label zu, aus dem sich die Entitäten ergeben
(Sequenzlabeling)~\cite{devlin2019bert,he2021debertav3}.
Large Language Models (LLMs) sind deutlich größere, generative
Sprachmodelle, die ein breites Spektrum von Aufgaben mit demselben Modell
bearbeiten~\cite{brown2020language,zhao2023survey}.
Sie lösen NER, indem sie die erkannten Entitäten als strukturierten Text
ausgeben~\cite{zhan2026generativeNER,semin2026spanLabeling}.
Im einfachsten Fall erhält das Modell dazu lediglich eine
Aufgabenbeschreibung im Prompt, ohne auf annotierten Daten trainiert zu
werden (\emph{Zero-Shot})~\cite{brown2020language}; alternativ wird es mit
annotierten Beispielen an die Aufgabe angepasst~\cite{zhan2026generativeNER}.

Die beiden Modellklassen unterscheiden sich nicht nur in ihrer
Arbeitsweise, sondern auch im Ressourcenbedarf.
Encoder wie DeBERTa-v3 umfassen einige hundert Millionen
Parameter~\cite{he2021debertav3}, aktuelle LLMs dagegen Milliarden bis
Hunderte Milliarden~\cite{qwen2026qwen35blog}.
Zudem erzeugen LLMs ihre Ausgabe Token für Token, sodass die Inferenzzeit
mit der Länge der Ausgabe wächst~\cite{lu2024padellm}, während ein Encoder
die Labels eines Satzes in einem einzigen Vorwärtsdurchlauf
bestimmt~\cite{devlin2019bert}.
Auch die Anpassung ist bei LLMs aufwendiger: Ein vollständiges
Fine-Tuning aller Gewichte erfordert zusätzlich Speicher für Gradienten
und Optimiererzustände~\cite{hu2021lora}.
Parameter-Efficient Fine-Tuning (PEFT) umfasst Verfahren, die stattdessen
nur eine kleine Zahl zusätzlicher Parameter trainieren~\cite{han2024parameter};
mit QLoRA lassen sich so auch Modelle mit Dutzenden Milliarden Parametern
auf einer einzelnen GPU anpassen~\cite{dettmers2023qlora}.
Für die praktische Modellwahl stellt sich damit die Frage, wann der
zusätzliche Aufwand eines LLM durch eine bessere Erkennungsleistung
gerechtfertigt ist und wann ein kompakter Encoder die effizientere Wahl
darstellt.

\section{Problemstellung, Forschungsfragen und Beitrag}
\label{sec:problem_ziel}

Mehrere Arbeiten vergleichen generative LLMs bereits mit Encoder-basierten
NER-Verfahren.
Ohne aufgabenspezifische Anpassung liegen LLMs dabei deutlich zurück:
In der Vergleichsstudie von \textcite{keraghel2024recent} erreicht GPT-4o
mit Few-Shot-Prompting auf CoNLL-2003 einen Makro-F$_1$ von 65{,}17, ein
vollständig angepasstes DeBERTa-v3-large dagegen 90{,}98.
Mit LoRA angepasste LLMs schließen diesen Abstand.
\textcite{zhan2026generativeNER} berichten für LLaMA3.1-8B im besten
Ausgabeformat einen F$_1$ von 93{,}85 auf CoNLL-2003 gegenüber 93{,}04 für
das Encoder-Verfahren BERT-MRC und zeigen zugleich, dass das Ausgabeformat
die Erkennungsleistung stark beeinflusst.
\textcite{obeidat2025llm} vergleichen auf fünf biomedizinischen
Datensätzen Encoder mit vollständigem Fine-Tuning und mittels QLoRA
angepasste LLMs mit 7 bis 8~Mrd. Parametern.
Die LLMs liegen auf vier Datensätzen um 2 bis 8 F$_1$-Punkte vorn,
benötigen je Testbeispiel aber das 40- bis 220-Fache der Inferenzzeit des
besten Encoders.

In der hier gesichteten Literatur lassen diese Vergleiche vier Fragen
offen, die Abschnitt~\ref{sec:research_gap} ausführlich herleitet.
Erstens verwenden sie entweder allgemeine Datensätze mit wenigen groben
Typen oder domänenspezifische Korpora; ein allgemeiner, feingranularer
Datensatz wurde für den direkten Vergleich nicht herangezogen.
Zweitens stammen die Encoder-Referenzwerte teilweise aus anderen
Publikationen statt aus derselben Pipeline.
Drittens werden Erkennungsleistung und Ressourcenbedarf nur teilweise
gemeinsam erfasst: \textcite{obeidat2025llm} messen Inferenzzeiten, aber
keinen Speicherbedarf, \textcite{zhan2026generativeNER} berichten nur
Trainingskosten.
Viertens stammen die untersuchten LLMs aus einer einzigen Größenklasse
zwischen 3 und 9~Mrd. Parametern, sodass offen bleibt, wie deutlich
kleinere oder größere Modelle im Vergleich zu Encodern abschneiden.

Die vorliegende Arbeit setzt an diesen Punkten an.
Sie vergleicht auf dem englischen Teil von MultiNERD mit seinen
15 Entitätstypen~\cite{tedeschi2022multinerd} drei Regime:
DeBERTa-v3-base und DeBERTa-v3-large mit vollständigem
Fine-Tuning~\cite{he2021debertav3}, Qwen3.5-0.8B, Qwen3.5-4B und
Qwen3.5-27B~\cite{qwen2026qwen35blog} im Zero-Shot-Einsatz sowie dieselben
Qwen3.5-Modelle nach einer Anpassung mit QLoRA.
Das Zero-Shot-Regime zeigt, was die veröffentlichten Modelle ohne
Trainingsaufwand leisten, und dient als Bezugspunkt für den Nutzen der
Anpassung.
QLoRA trainiert Low-Rank-Adapter auf einem in 4~Bit quantisierten,
eingefrorenen Basismodell~\cite{hu2021lora,dettmers2023qlora} und macht
so auch die Anpassung des 27B-Modells auf einer einzelnen GPU möglich.
Alle Konfigurationen verwenden dieselben offiziellen Datensplits und werden
mit derselben strikten Metrik auf Entitätsebene bewertet; die generierten
Ausgaben der LLMs werden dazu in dieselbe Label-Repräsentation überführt
wie die Vorhersagen der Encoder.
Neben der Erkennungsleistung werden die Trainingsdauer, die
Inferenzlatenz und der GPU-Spitzenspeicher während der Inferenz erfasst.

Daraus ergeben sich drei Forschungsfragen (FF1--FF3):
\begin{description}
  \item[FF1 (Erkennungsleistung):] Wie unterscheiden sich
  DeBERTa-v3-Encoder mit vollständigem Fine-Tuning und Qwen3.5-LLMs im
  Zero-Shot- und im QLoRA-Regime in ihrer Erkennungsleistung auf dem
  englischen MultiNERD-Testsplit, und wie zuverlässig halten die LLMs das
  vorgegebene Ausgabeformat ein (Formtreue)?
  \item[FF2 (Modellgröße und Ressourcen):] Wie verändern sich
  Erkennungsleistung, Trainingsdauer, Inferenzlatenz und
  GPU-Spitzenspeicher der Inferenz mit der Modellgröße innerhalb beider
  Modellfamilien?
  \item[FF3 (Empfehlungen):] Welche Konfigurationen bieten unter gegebenen
  Speicher- und Latenzbeschränkungen das günstigste Verhältnis von
  Erkennungsleistung und Ressourcenbedarf, und unter welchen Bedingungen
  rechtfertigt ein LLM den Mehraufwand gegenüber einem Encoder?
\end{description}

Die Arbeit leistet damit folgende Beiträge:
\begin{itemize}
  \item einen Vergleich von acht Konfigurationen beider Modellklassen in
  einer gemeinsamen Pipeline mit eigener Encoder-Referenz statt
  übernommener Werte; die fünf trainierten Konfigurationen werden mit je
  drei Seeds ausgewertet, sodass sich insgesamt 18 Testevaluationen
  ergeben;
  \item eine Größenstaffelung über zwei Modellfamilien von rund 184~Mio.
  bis 27~Mrd. Parametern mit gemeinsamer Erfassung von
  Erkennungsleistung und Ressourcenbedarf;
  \item ein JSON-Ausgabeformat, das Entitäten über Wortindizes adressiert
  und ohne Constrained Decoding erzeugt wird, sodass die Formtreue der
  LLM-Ausgaben eigens gemessen werden kann~\cite{semin2026spanLabeling};
  \item einen ergänzenden Transfertest auf CrossNER, der prüft, wie sich
  die jeweils größten auf MultiNERD angepassten Modelle beider Familien
  auf Wikipedia-Texte aus fünf anderen Domänen übertragen
  lassen~\cite{liu2021crossner};
  \item eine Absicherung der Befunde durch gepaarte
  Bootstrap-Konfidenzintervalle über die Testsätze und durch zwei
  Sensitivitätsanalysen zum Ausgabelimit der LLM-Inferenz, die den
  wichtigsten Störfaktor des Vergleichs messen statt schätzen.
\end{itemize}
Der Transfertest verwendet ausschließlich die Testsplits von CrossNER,
ohne weiteres Training, und wertet nur die gemeinsamen Grobtypen Person,
Organisation und Ort aus.
Er dient der Einordnung der Hauptergebnisse, nicht dem Vergleich mit
Bestwerten der Literatur.

\section{Abgrenzung der Arbeit}
\label{sec:abgrenzung}

Der angemeldete Titel der Arbeit nennt BERT-basierte Modelle und LoRA.
Auf der Encoder-Seite werden mit DeBERTa-v3-base und DeBERTa-v3-large
zwei Weiterentwicklungen von BERT untersucht, die sich in vergleichenden
NER-Studien als besonders leistungsfähig erwiesen haben
(Abschnitt~\ref{sec:encoder_ner}); auf der LLM-Seite wird LoRA in seiner
speicheroptimierten Form QLoRA eingesetzt, die das Training des
27B-Modells auf einer einzelnen GPU erst möglich macht
(Abschnitt~\ref{sec:lora}), und das Zero-Shot-Regime kommt als
Bezugspunkt ohne Training hinzu.
Im Folgenden werden die Verfahren daher mit diesen präziseren
Bezeichnungen benannt.

Die Untersuchung beschränkt sich auf NER in englischsprachigen Texten.
Hauptdatenbasis sind die offiziellen Splits des englischen Teils von
MultiNERD mit Texten aus Wikipedia und Wikinews~\cite{tedeschi2022multinerd}.
Da diese Splits automatisch annotiert wurden, messen die berichteten Werte
die Übereinstimmung mit einer automatisch erzeugten Referenz
(Abschnitte~\ref{sec:datasets_corpora} und~\ref{sec:datenauswahl}).
CrossNER wird ausschließlich als Transfertest auf den Testsplits und für
die Typen Person, Organisation und Ort ausgewertet~\cite{liu2021crossner}.
CoNLL-2003 und OntoNotes dienen nur der Einordnung verwandter Arbeiten und
werden nicht experimentell untersucht.
Eine Übertragung der Schlussfolgerungen auf andere Sprachen oder
Anwendungsdomänen wird nicht abgeleitet.

Auf der Encoder-Seite werden DeBERTa-v3-base und DeBERTa-v3-large
eingesetzt, die sich in vergleichenden NER-Untersuchungen als besonders
leistungsfähig erwiesen haben~\cite{keraghel2024recent}.
Sie werden um einen linearen Klassifikationskopf ergänzt und
ausschließlich mit vollständigem Fine-Tuning angepasst, bei dem
Encoder-Gewichte und Klassifikationskopf gemeinsam trainiert
werden~\cite{devlin2019bert}; ein Zero-Shot-Einsatz ist ohne trainierten
Klassifikationskopf nicht möglich.
Auf der LLM-Seite werden drei Größen der Qwen3.5-Familie untersucht, die
einen Vergleich innerhalb einer Modellfamilie erlauben~\cite{qwen2026qwen35blog}.
Sie werden ohne Training im Zero-Shot-Regime sowie nach einer Anpassung
mit QLoRA eingesetzt, bei der nur die Low-Rank-Adapter trainiert
werden~\cite{hu2021lora,dettmers2023qlora}.
Nicht untersucht werden ein vollständiges Fine-Tuning der LLMs,
alternative PEFT-Verfahren, weitere Modellfamilien, Few-Shot-Prompting
mit Demonstrationsbeispielen, Constrained Decoding und der Thinking-Modus
der Qwen3.5-Modelle; andere Modellfamilien und Verfahren werden nur im
Literaturüberblick berücksichtigt.

Die Regime unterscheiden sich in mehreren Faktoren zugleich, darunter
Vortraining, Architektur, Modellgröße, Quantisierung und
Anpassungsverfahren; zudem liefen die Konfigurationen auf
unterschiedlichen GPU-Typen (Abschnitt~\ref{sec:validitaet}).
Die Ergebnisse sind daher als Vergleich konkreter Modell- und
Anpassungskonfigurationen zu lesen; isolierte kausale Effekte einzelner
Faktoren werden mit diesem Versuchsdesign nicht bestimmt.

Die Erkennungsleistung wird mit Precision, Recall und F$_1$-Score auf
Entitätsebene gemessen; eine Entität gilt nur dann als korrekt erkannt,
wenn ihre Grenzen und ihr Typ mit der Referenz
übereinstimmen~\cite{manning2008introduction,sang2003conll}.
Zentrale Kennzahl ist der Mikro-F$_1$-Score.
Für die LLMs wird zusätzlich die Formtreue ihrer Ausgaben erfasst
(Abschnitt~\ref{sec:robustheit-ausgaben}).
Der Ressourcenbedarf umfasst die Trainingsdauer, die Inferenzlatenz bei
einzeln verarbeiteten Sätzen und den GPU-Spitzenspeicher während der
Inferenz (Abschnitt~\ref{sec:effizienzmetriken}).
Nicht erfasst werden der Speicherbedarf während des Trainings, ein durch
Batching optimierter Durchsatz sowie Energieverbrauch und Kosten.
Eine mögliche Kontamination der Vortrainingsdaten der LLMs mit
MultiNERD-Texten wird diskutiert, aber nicht empirisch geprüft
(Abschnitt~\ref{sec:research_gap}).
Interpretierbarkeit, Fairness und Robustheit gegenüber adversarialen
Eingaben sind ebenso wenig Gegenstand der Arbeit wie die Erhebung und
Annotation eines neuen Datensatzes.

\section{Aufbau der Arbeit}
\label{sec:aufbau_arbeit}

Die weiteren Kapitel führen von den Grundlagen und dem Forschungsstand
über die Methodik zu den Ergebnissen und Schlussfolgerungen.

Kapitel~\ref{chap:grundlagen} führt in NLP und NER ein und beschreibt die
Transformer-Architektur, Encoder-Modelle für NER, generative LLMs sowie
parameter-effizientes Fine-Tuning mit LoRA und QLoRA.
Nach einer Einordnung von Skalierungsgesetzen stellt es die verwendeten
Datensätze, darunter MultiNERD und CrossNER, sowie die
Evaluationsmetriken vor.

Kapitel~\ref{chap:verwandte_arbeiten} stellt Encoder-basierte und
LLM-basierte NER-Ansätze einschließlich verschiedener Ausgabeformate vor.
Es behandelt Arbeiten, die mit PEFT angepasste LLMs direkt mit Encodern
vergleichen, fasst Befunde zum Einfluss der Modellgröße zusammen und
leitet daraus die Forschungslücke ab.

Kapitel~\ref{chap:methodik} beschreibt die drei Vergleichsregime, die
Auswahl von Modellen und Daten sowie die Encoder- und die LLM-Pipeline
einschließlich Prompt-Design, QLoRA-Konfiguration und Nachverarbeitung der
generierten Ausgaben.
Darauf folgen das Evaluationskonzept für Erkennungsleistung,
Ressourcenbedarf und Formtreue, das Versuchsdesign des CrossNER-Transfertests
und eine Zusammenfassung der bekannten Abweichungen vom idealisierten
Versuchsdesign.

Kapitel~\ref{chap:ergebnisse} präsentiert die Erkennungsleistung aller
Konfigurationen insgesamt und je Entitätstyp sowie die Formtreue der
LLM-Ausgaben, untersucht den Einfluss der Modellgröße und den
Ressourcenbedarf, analysiert die Fehlerprofile der Modellklassen,
berichtet die Ergebnisse des CrossNER-Transfertests und diskutiert die
Befunde.

Kapitel~\ref{chap:fazit} beantwortet die Forschungsfragen, erläutert die
Grenzen der Untersuchung und skizziert weiterführende Arbeiten.
Der Anhang dokumentiert die Reproduzierbarkeit der Läufe mit
Konfigurationen, Codeständen, Modellrevisionen und nachrangigen
Einstellungen sowie den vollständigen Zero-Shot-Prompt.
